\section{LL gramatika}
\label{app:gramatika}

$G = (N, T, P, \texttt{<program>})$, kde
\begin{itemize}\raggedright
  \item $N = \{$\texttt{<program>}, \texttt{<prolog>}, \texttt{<funcs>}, \texttt{<func\_def>}, \texttt{<main\_opt>}, \texttt{<nl\_opt>}, \texttt{<func\_body>}, \texttt{<params>}, \texttt{<params\_tail>}, \texttt{<param>}, \texttt{<type\_opt>}, \texttt{<type>}, \texttt{<block>}, \texttt{<statements>}, \texttt{<statements\_tail>}, \texttt{<statement>}, \texttt{<stmt\_sep>}, \texttt{<simple\_stmt>}, \texttt{<var\_def>}, \texttt{<var\_kind>}, \texttt{<return>}, \texttt{<ret\_value>}, \texttt{<if>}, \texttt{<while>}, \texttt{<call>}, \texttt{<args>}, \texttt{<args\_tail>}$\}$,
  \item $T = \{$\texttt{id}, \texttt{import}, \texttt{def}, \texttt{val}, \texttt{var}, \texttt{if}, \texttt{then}, \texttt{else}, \texttt{while}, \texttt{do}, \texttt{return}, \texttt{Int}, \texttt{Double}, \texttt{String}, \texttt{Unit}, \texttt{(}, \texttt{)}, \texttt{,}, \texttt{:}, \texttt{;}, \texttt{.}, \texttt{*}, \texttt{=}, \texttt{@}, \texttt{NEWLINE}, \texttt{INDENT}, \texttt{DEDENT}, \texttt{EOF}, \texttt{expr}$\}$,
  \item $P$ obsahuje pravidla 1--48:
\end{itemize}

\texttt{expr} je brán jako terminál, protože jej analyzuje modul \texttt{expression\_parser.c/h} metodou precedenční analýzy. Kvůli implementaci rozšíření FUNEXP je \texttt{<call>} součástí \texttt{expr}. 

Pravidla 44--48 jsou samostatná podgramatika, zpracovávaná rekurzivním sestupem, se startovním symbolem \texttt{<call>}, který je nedosažitelný z \texttt{<program>}. Volá je precedenční analýza, když je na vstupu dvojice \texttt{id} \texttt{(}. Každý argument je zase \texttt{expr}, což dovoluje modulu analyzovat i vnořená volání \texttt{f(g(x))}.

\subsection*{Program a funkce}
{\small
\begin{tabular}{@{}r@{\hspace{0.6em}}p{9.5em}@{\hspace{0.4em}}c@{\hspace{0.4em}}l@{}}
1: & \texttt{<program>} & $\rightarrow$ & \texttt{<prolog>} \texttt{<funcs>} \texttt{EOF} \\
2: & \texttt{<prolog>} & $\rightarrow$ & \texttt{import} \texttt{id} \texttt{.} \texttt{*} \texttt{NEWLINE} \\
3: & \texttt{<funcs>} & $\rightarrow$ & \texttt{<func\_def>} \texttt{<funcs>} \\
4: & \texttt{<funcs>} & $\rightarrow$ & $\varepsilon$ \\
5: & \texttt{<func\_def>} & $\rightarrow$ & \texttt{<main\_opt>} \texttt{def} \texttt{id} \texttt{(} \texttt{<params>} \texttt{)} \texttt{<type\_opt>} \texttt{=} \texttt{<func\_body>} \\
6: & \texttt{<main\_opt>} & $\rightarrow$ & \texttt{@} \texttt{id} \texttt{<nl\_opt>} \\
7: & \texttt{<main\_opt>} & $\rightarrow$ & $\varepsilon$ \\
8: & \texttt{<nl\_opt>} & $\rightarrow$ & \texttt{NEWLINE} \\
9: & \texttt{<nl\_opt>} & $\rightarrow$ & $\varepsilon$ \\
10: & \texttt{<func\_body>} & $\rightarrow$ & \texttt{expr} \texttt{NEWLINE} \\
11: & \texttt{<func\_body>} & $\rightarrow$ & \texttt{<block>} \\
\end{tabular}}

\subsection*{Parametry a typy}
{\small
\begin{tabular}{@{}r@{\hspace{0.6em}}p{9.5em}@{\hspace{0.4em}}c@{\hspace{0.4em}}l@{}}
12: & \texttt{<params>} & $\rightarrow$ & \texttt{<param>} \texttt{<params\_tail>} \\
13: & \texttt{<params>} & $\rightarrow$ & $\varepsilon$ \\
14: & \texttt{<params\_tail>} & $\rightarrow$ & \texttt{,} \texttt{<param>} \texttt{<params\_tail>} \\
15: & \texttt{<params\_tail>} & $\rightarrow$ & $\varepsilon$ \\
16: & \texttt{<param>} & $\rightarrow$ & \texttt{id} \texttt{:} \texttt{<type>} \\
17: & \texttt{<type\_opt>} & $\rightarrow$ & \texttt{:} \texttt{<type>} \\
18: & \texttt{<type\_opt>} & $\rightarrow$ & $\varepsilon$ \\
19: & \texttt{<type>} & $\rightarrow$ & \texttt{Int} \\
20: & \texttt{<type>} & $\rightarrow$ & \texttt{Double} \\
21: & \texttt{<type>} & $\rightarrow$ & \texttt{String} \\
22: & \texttt{<type>} & $\rightarrow$ & \texttt{Unit} \\
\end{tabular}}

\subsection*{Bloky a příkazy}
{\small
\begin{tabular}{@{}r@{\hspace{0.6em}}p{9.5em}@{\hspace{0.4em}}c@{\hspace{0.4em}}l@{}}
23: & \texttt{<block>} & $\rightarrow$ & \texttt{NEWLINE} \texttt{INDENT} \texttt{<statements>} \texttt{DEDENT} \\
24: & \texttt{<statements>} & $\rightarrow$ & \texttt{<statement>} \texttt{<statements\_tail>} \\
25: & \texttt{<statements\_tail>} & $\rightarrow$ & \texttt{<statement>} \texttt{<statements\_tail>} \\
26: & \texttt{<statements\_tail>} & $\rightarrow$ & $\varepsilon$ \\
27: & \texttt{<statement>} & $\rightarrow$ & \texttt{<simple\_stmt>} \texttt{<stmt\_sep>} \\
28: & \texttt{<statement>} & $\rightarrow$ & \texttt{<if>} \\
29: & \texttt{<statement>} & $\rightarrow$ & \texttt{<while>} \\
30: & \texttt{<stmt\_sep>} & $\rightarrow$ & \texttt{NEWLINE} \\
31: & \texttt{<stmt\_sep>} & $\rightarrow$ & \texttt{;} \texttt{<nl\_opt>} \\
32: & \texttt{<simple\_stmt>} & $\rightarrow$ & \texttt{<var\_def>} \\
33: & \texttt{<simple\_stmt>} & $\rightarrow$ & \texttt{<return>} \\
34: & \texttt{<simple\_stmt>} & $\rightarrow$ & \texttt{id} \texttt{=} \texttt{expr} \\
35: & \texttt{<simple\_stmt>} & $\rightarrow$ & \texttt{expr} \\
36: & \texttt{<var\_def>} & $\rightarrow$ & \texttt{<var\_kind>} \texttt{id} \texttt{<type\_opt>} \texttt{=} \texttt{expr} \\
37: & \texttt{<var\_kind>} & $\rightarrow$ & \texttt{val} \\
38: & \texttt{<var\_kind>} & $\rightarrow$ & \texttt{var} \\
39: & \texttt{<return>} & $\rightarrow$ & \texttt{return} \texttt{<ret\_value>} \\
40: & \texttt{<ret\_value>} & $\rightarrow$ & \texttt{expr} \\
41: & \texttt{<ret\_value>} & $\rightarrow$ & $\varepsilon$ \\
42: & \texttt{<if>} & $\rightarrow$ & \texttt{if} \texttt{expr} \texttt{then} \texttt{<block>} \texttt{else} \texttt{<block>} \\
43: & \texttt{<while>} & $\rightarrow$ & \texttt{while} \texttt{expr} \texttt{do} \texttt{<block>} \\
\end{tabular}}

\subsection*{Volání funkcí (součást expr v precedenční analýze)}
{\small
\begin{tabular}{@{}r@{\hspace{0.6em}}p{9.5em}@{\hspace{0.4em}}c@{\hspace{0.4em}}l@{}}
44: & \texttt{<call>} & $\rightarrow$ & \texttt{id} \texttt{(} \texttt{<args>} \texttt{)} \\
45: & \texttt{<args>} & $\rightarrow$ & \texttt{expr} \texttt{<args\_tail>} \\
46: & \texttt{<args>} & $\rightarrow$ & $\varepsilon$ \\
47: & \texttt{<args\_tail>} & $\rightarrow$ & \texttt{,} \texttt{expr} \texttt{<args\_tail>} \\
48: & \texttt{<args\_tail>} & $\rightarrow$ & $\varepsilon$ \\
\end{tabular}}


\subsection*{Poznámky k pravidlům}
\begin{itemize}
  \item pravidlo 2: id musí být \texttt{ifj26}, jinak chyba 2
  \item pravidlo 6: id musí být \texttt{main}, jinak chyba 2
  \item pravidlo 10: jednořádková funkce
  \item pravidlo 11: víceřádková funkce
  \item pravidlo 34: LL(2): \texttt{id} a za ním \texttt{=}
\end{itemize}
